\section{Local unitary densities and the exact mass ratio}
\label{md-section}

We verify the local densities and the mass normalization used in
Section~\ref{as-section}. The finite-field calculations and all
normalization changes are proved below, including residue
characteristic $2$. The global mass formula is used as an explicitly
stated classical theorem. This is sufficient to check the source's
mass-ratio assertion; it is not a new proof of the general mass formula.

\subsection{Counting primitive frames over a finite field}

Let $k=\mathbb F_q$, where $q$ is a prime power, and let
$k'=\mathbb F_{q^2}$ with involution $a\mapsto\overline a=a^q$.
On $(k')^n$ put $h(x,y)=\overline x^{\,t}y$.
For a matrix $G=\overline G^{\,t}\in M_m(k')$ of rank $m-c$, let
$\mathcal E_-(n,G)$ be the set of matrices $X\in M_{n,m}(k')$
of column rank $m$ satisfying $\overline X^{\,t}X=G$.
In the split case, let $G\in M_m(k)$ have rank $m-c$, and let
$\mathcal E_+(n,G)$ consist of pairs $(X,Y)$ of matrices in
$M_{n,m}(k)$, both of column rank $m$, satisfying $X^tY=G$.
In both cases $0\leq c\leq m$.
For $\eta\in\{1,-1\}$ define
\begin{equation}\label{md-density}
 d_\eta(n,m,c;q)=
 \begin{cases}
 \displaystyle\prod_{j=n-m-c+1}^{n}(1-\eta^j q^{-j}),&m+c\leq n,\\
 0,&m+c>n.
 \end{cases}
\end{equation}
An empty product is $1$. Equivalently, one may start the product at
$\max(0,n-m-c+1)$: the extra cases then contain the zero factor
$1-\eta^0q^0$.

\begin{lemma}\label{md-unitary-order}
For $n\geq0$,
\[
 \#U_n(k)=q^{n(n-1)/2}\prod_{j=1}^n(q^j-(-1)^j).
\]
For $n\geq1$, the number of nonzero isotropic vectors in $(k')^n$ is
\[
                  (q^n-(-1)^n)(q^{n-1}-(-1)^{n-1}).
\]
Every nondegenerate Hermitian form on $(k')^n$ is isometric to $h$.
These assertions hold also when $q$ is even.
\end{lemma}

\begin{proof}
The norm $N:k'^\times\to k^\times$, $N(a)=a^{q+1}$, is
surjective and every fiber has $q+1$ elements: the multiplicative
group of a finite field is cyclic. The trace $a+\overline a$ is
a nonzero $k$-linear map $k'\to k$, hence is also surjective.
It is nonzero even in characteristic $2$, since the fixed field
of the involution is exactly $k$.

A nonzero nondegenerate Hermitian form has a vector of nonzero
length. Otherwise, expanding $h(ax+y,ax+y)=0$ for all $a\in k'$
would give $\operatorname{Tr}(\overline a h(x,y))=0$ for all $a$.
The trace pairing is nondegenerate, so the form would be zero.
Normalize a vector of nonzero length using the norm map, take its
nondegenerate orthogonal complement, and induct. This proves the
last assertion and extends every orthonormal frame to a basis.

Let $A_n$ count vectors of length $0$ and let $B_n$ count vectors
of length a prescribed nonzero element of $k$. The latter count
is independent of that element by norm surjectivity. We have
$A_0=1$, $B_0=0$, and, by separating the last coordinate,
\[
 A_n=A_{n-1}+(q^2-1)B_{n-1},\qquad
 A_n+(q-1)B_n=q^{2n}.
\]
Equivalently $A_n-B_n=-q(A_{n-1}-B_{n-1})=(-q)^n$.
Solving these two equations gives
\[
 B_n=q^{n-1}(q^n-(-1)^n),\qquad
 A_n=q^{2n-1}+(q-1)(-1)^nq^{n-1}.
\]
Thus $A_n-1$ is the second expression in the lemma.
Choosing the first vector of an orthonormal basis in $B_n$ ways
and then a basis of its orthogonal complement gives
$\#U_n(k)=B_n\#U_{n-1}(k)$, proving the product formula.
\end{proof}

\begin{theorem}\label{md-frame-count}
For the two sets of primitive frames defined above,
\[
 \#\mathcal E_-(n,G)=q^{2nm-m^2}d_{-1}(n,m,c;q),\qquad
 \#\mathcal E_+(n,G)=q^{2nm-m^2}d_1(n,m,c;q).
\]
The same formulas hold for any nondegenerate target Hermitian
form in the first case and any perfect target bilinear pairing
in the second case.
\end{theorem}

\begin{proof}
Changes of target bases remove the target form or pairing.
If $m>n$ both sets are empty, and the displayed products are zero.
The assertion is immediate if $m=0$. We may therefore assume
$1\leq m\leq n$.

First consider the Hermitian case. A basis of the radical and
Lemma~\ref{md-unitary-order} put $G$ into the form
$\operatorname{diag}(I_r,0_c)$, where $r=m-c$.
Choose the first $r$ columns as an orthonormal frame. There are
$\#U_n(k)/\#U_{n-r}(k)$ choices. The last $c$ columns are a
linearly independent ordered isotropic frame in the orthogonal
complement of dimension $N=n-r$; all their pairings must vanish.

Let $F(N,c)$ count such ordered isotropic frames. Their span is
contained in its own orthogonal complement, so $2c\leq N$ is
necessary. If $c\geq1$, choose the first nonzero isotropic vector
$v$. The quotient $v^\perp/k'v$ is nondegenerate of dimension
$N-2$. To see this directly, choose $w$ with $h(v,w)=1$ and
replace $w$ by $w+av$ so that it is isotropic; the required
choice is possible because the trace is surjective. The span
of $v,w$ is then a nondegenerate plane.
Every isotropic $(c-1)$-frame in the quotient has
$q^{2(c-1)}$ lifts to $v^\perp$, all of which form with $v$
an independent isotropic frame. Therefore
\[
 F(N,c)=(q^N-(-1)^N)(q^{N-1}-(-1)^{N-1})
                         q^{2(c-1)}F(N-2,c-1).
\]
Together with $F(N,0)=1$, this proves
\[
 F(N,c)=q^{c(c-1)}
           \prod_{j=N-2c+1}^{N}(q^j-(-1)^j)
                  \quad(2c\leq N).
\]
The necessary condition $2c\leq n-r$ is precisely $m+c\leq n$.
Multiplying by the orthonormal-frame count makes the two products
adjacent, with indices $n-m-c+1,\ldots,n$. The total exponent of
$q$ after writing each factor as $q^j(1-(-1)^jq^{-j})$ is
\[
 \frac{n(n-1)-(n-r)(n-r-1)}2+c(c-1)
             +\sum_{j=n-m-c+1}^{n}j
                       =2nm-m^2.
\]
This proves the Hermitian formula, including the impossible cases.

For the split case, the number of possible matrices $X$ is
$\prod_{i=0}^{m-1}(q^n-q^i)$. For a fixed $X$, a change of
dual target bases reduces the equation to saying that the first
$m$ rows of $Y$ are $G$. Write the remaining $n-m$ rows as $Z$.
The stacked matrix $Y$ is injective precisely when $Z$ is
injective on $\ker G$, which has dimension $c$. After choosing
a complement of this kernel, the number of choices is
\[
           q^{(n-m)(m-c)}\prod_{i=0}^{c-1}(q^{n-m}-q^i).
\]
It is zero if $c>n-m$. Otherwise, multiply by the count of $X$
and extract powers of $q$. The remaining product has exactly
the indices $n-m-c+1,\ldots,n$, and the extracted exponent is
$2nm-m^2$. This proves the split formula.
\end{proof}

\subsection{The sign of the primitive density error}
\label{sec:primitive-local-density}

For a self-dual unramified local Hermitian lattice of rank $N$, let
$\delta(A)$ be the primitive residue-frame density of an integral Gram
matrix $A$ of size $m$ and residue corank $c$, with $m+c\le N$.
Primitive frame counting, with denominator $q^{2Nm-m^2}$, gives
\begin{equation}
 \delta(A)=\prod_{j=d}^{N}(1-\eta^j q^{-j}),
 \qquad d=N-m-c+1,
 \label{eq:primitive-density-product}
\end{equation}
where $\eta=1$ at a split place and $\eta=-1$ at an inert place.
This is the reconstructed primitive-frame count \cite{reconstruction}.
The formula includes singular residue matrices and characteristic two.
At a designated split level the original tuple
$(f_a,\varpi^{-1}(f'_a-f_a))_a$ must be retained.

In the split case all factors lie in $(0,1)$, so
$0\le1-\delta(A)\le q^{-d}/(1-q^{-1})$.
At an inert place the sharper signed bounds are
\[
 \begin{cases}
  1-q^{-d}\le\delta(A)\le1,&d\text{ even},\\
  1\le\delta(A)\le1+q^{-d},&d\text{ odd}.
 \end{cases}
\]
To prove them, pair consecutive factors. A pair beginning at an odd
index is $(1+t)(1-t/q)\ge1$, and one beginning at an even index is
$(1-t)(1+t/q)\le1$. An unmatched last factor has the same sign of
deviation as the first. Removing the first factor gives the opposite
endpoint bound.

For the original even parameters $m=2r+2$ and $N=k+m-1$, put
$d_0=k-m\ge4$. Then $d=k-c$ and
$|\delta-1|\le2q^{-d_0}$. Positive errors occur only at inert places
with odd corank, hence $c\le m-1$ and
\[
 0<\delta\le\exp(q^{-(d_0+1)}),\qquad
 \log\prod_v\delta_v
 \le f\,2^{-(d_0+1)}\left(1+\frac2{d_0}\right)
\]
over a number field of degree $f$. The product converges absolutely;
the last bound follows by summing first over places above each rational
prime, then bounding the tail by an integral. This sharpens the
finite-place coefficient without supplying the archimedean orbital
estimates or changing the restrictions of the period maps.

\subsection{Why equivalent local data have the same density}
\label{sec:scalar-pair-local-density}

Let $F$ be a nonarchimedean local field of characteristic zero,
$E/F$ an unramified quadratic field or $E=F\times F$, and $O$ its
maximal integral algebra. In nondegenerate Hermitian $E$-spaces of the
same rank $N$, let $\Lambda_i$ be self-dual lattices and
$(f_i,f'_i)$ pairs of frames of length $r+1$. In each frame the first
$r$ vectors are totally isotropic and perpendicular to the last.
Assume the combined raw spans are nondegenerate. Let $M_i$ be generated
by $(f_i,f'_i)$ away from the designated split level and by
\[
       \left(f_i,\frac{f'_i-f_i}{\varpi_s}\right)
\]
at that level. Assume the displayed tuples are linearly independent
and primitive in $\Lambda_i$. The generator $\varpi_s$ is prescribed.
Put $P_i=EM_i$, $W_i=P_i^\perp$ and $K_i=\Lambda_i\cap W_i$.
Let $X_i,X'_i$ denote the spans of the first $r$ frame vectors.

Assume the two quotient spaces $X_i^\perp/X_i$ and
$(X'_i)^\perp/X'_i$ are identified isometrically, with the following
compatibilities. Each vector of $W_1$ has a representative in $W_2$
with both transported quotient classes; the integral quotient-lattice
images are preserved; and, at the split level, the relation represented
by integral pairs $u,u'$ with $u-u'\in\varpi_s\Lambda_i$ is preserved
in both directions. These are among the original equivalence
conditions.

\begin{proposition}
The lattices $K_1,K_2$ are isometric. The discriminant modules of
$M_1,M_2$ are isomorphic, and their original primitive local densities
are equal.
\end{proposition}

\begin{proof}
For the original data, writing the last frame norms as $-T_i,-T'_i$,
the raw frame Gram determinant is
\[
 (-1)^{r+1}
 \left(\operatorname{Nm}\det Z_i
       -T_iT'_i\operatorname{Nm}\det Z_i^\circ\right)\ne0,
\]
where $Z_i=(h(f'_{i,b},f_{i,a}))_{a,b}$ and $Z_i^\circ$ deletes its
last row and column. The rank-two determinant formula gives this
identity, and the original hypotheses give its nonvanishing.
Thus the assumed nondegeneracy is satisfied. Since
$W_i\cap X_i=0$, the transported representative is unique and defines
an isometry $\alpha:W_1\to W_2$. It is onto by equal dimensions,
componentwise in the split algebra.

For $w\in K_1$, put $w_2=\alpha(w)$. Away from the designated level,
the integral quotient images give
$u=w_2+\sum a_jf_{2,j}$ and
$u'=w_2+\sum b_jf'_{2,j}$ in $\Lambda_2$.
Their difference belongs to $\Lambda_2\cap P_2=M_2$.
Primitivity of the raw tuple forces the coefficients to be integral,
hence $w_2$ is integral.

At the designated level, the congruence relation instead gives these
representatives with $u-u'\in\varpi_s\Lambda_2$. Extend their
first-$r$ coefficient lists by zero in the last position. Then
\[
 \frac{u-u'}{\varpi_s}
   =\sum_j\frac{a_j-b_j}{\varpi_s}f_{2,j}
     -\sum_j b_j\frac{f'_{2,j}-f_{2,j}}{\varpi_s}\in M_2.
\]
Primitivity of the modified tuple forces
$b_j,(a_j-b_j)/\varpi_s\in O$, so again $w_2$ is integral.
The inverse identifications prove $\alpha(K_1)=K_2$.

For a primitive nondegenerate sublattice $M$ of a self-dual lattice
$\Lambda$, its orthogonal intersection $K$ is also primitive.
Both are module direct summands. Restricting the self-duality pairing
to $M$, and then to $K$, gives
\[
 M^\#/M\ \simeq\ \Lambda/(M+K)\ \simeq\ K^\#/K.
\]
The isometry of the $K_i$ therefore identifies the discriminant
modules of the $M_i$. Reducing these modules modulo a uniformizer
shows equality of the residue Gram coranks, in each split component.
The reconstructed primitive-frame formula depends only on this
corank and the common ranks and residue field \cite{reconstruction}.
It gives the required equality with the original normalizers.
\end{proof}

For the original ranks $m=2r+2$, $N=k+m-1$ and $k-m\ge4$, put
$q=\#k_F$. Together with $|\delta-1|\le2q^{-(k-2r-2)}$ and $2r+2\ge2$,
this proves the full original local density assertion with constant
one. It does not supply the remaining orbital or global norm estimates.


\subsection{Lifting primitive frames, including the dyadic case}

Let $A$ be a complete discrete valuation ring with uniformizer
$\varpi$ and residue field $k$ of size $q$. Let $B$ be either
the integers in the unramified quadratic extension of its fraction
field or $A\times A$ with its interchange involution. A Hermitian
matrix over $B$ is a matrix $G=\overline G^{\,t}$. In the split
case it is equivalent to one arbitrary matrix over $A$.
Let $\Phi\in M_n(B)$ be an invertible Hermitian matrix, and let
$G\in M_m(B)$ be Hermitian with reduction of corank $c$.
A primitive solution modulo $\varpi^a$ to
\[
                         \overline X^{\,t}\Phi X=G
\]
means one whose reduction has full column rank at both residue
factors, if there are two. Denote their number by $N_a$.

\begin{proposition}\label{md-lifting}
If $N_1\ne0$, then for every $a\geq1$,
\[
 N_a=q^{(a-1)(2nm-m^2)}N_1.
\]
The normalized primitive density is consequently
\[
 \lim_{a\to\infty}q^{-a(2nm-m^2)}N_a
                        =d_\eta(n,m,c;q),
\]
with $\eta=-1$ in the inert case and $\eta=1$ in the split case.
If $N_1=0$, all the counts and the density are zero.
\end{proposition}

\begin{proof}
Suppose first that the residue algebra is $k'$. A tentative lift
of a solution at level $a$ may be corrected by $\varpi^aH$.
Modulo $\varpi^{a+1}$ its Gram matrix changes by
\[
        \varpi^a(\overline X^{\,t}\Phi H+
                                      \overline H^{\,t}\Phi X).
\]
All matrices on the right can be reduced modulo $\varpi$.
Since $X$ is injective and $\Phi$ is invertible,
$\overline X^{\,t}\Phi$ has rank $m$. Thus
$T=\overline X^{\,t}\Phi H$ ranges over all of $M_m(k')$.
The map $T\mapsto T+\overline T^{\,t}$ is onto the
$k$-space of Hermitian matrices: choose upper off-diagonal
entries freely and choose each diagonal entry using the
surjective trace $k'\to k$. This argument uses no division by $2$.
The source space of corrections has $k$-dimension $2nm$ and the
target has dimension $m^2$, so every allowed Gram correction has
$q^{2nm-m^2}$ preimages.

In the split case write the two components as $X,Y$.
The linearized Gram map is
$\delta X^t\Phi Y+X^t\Phi\delta Y$. Already its second
term is onto $M_m(k)$, since $X$ has rank $m$.
The same dimension count gives the same number of lifts.
Iterating proves the first formula; division by the indicated
power of $q$ and Theorem~\ref{md-frame-count} prove the second.
Primitivity is preserved under lifting.
\end{proof}

In particular the unitary stabilizer density of a unimodular
rank-$n$ lattice is
\begin{equation}\label{md-stabilizer-density}
                 \beta_{n,v}=
                       \prod_{j=1}^n(1-\eta_v^j q_v^{-j}).
\end{equation}
This follows by taking $m=n$, $c=0$: a matrix with the prescribed
invertible Gram matrix is automatically invertible. The argument
does not require the local base field to be unramified over
$\mathbb Q_2$. It only requires that its quadratic algebra $B$
be unramified or split.

\begin{example}\label{md-degenerate-test}
For $q=2$, $n=2$, $m=1$, and zero Gram matrix, there are $9$
primitive Hermitian vectors and $3$ split pairs. Their normalized
densities are $9/8$ and $3/8$, respectively. For nonzero Gram
matrix the two counts are both $6$, and the two densities are
$3/4$. These values follow directly from the formulas above and
show why the corank $c$ must occur in the lower limit of the product.
Also, in rank $n=m=1$, a zero Gram matrix has no primitive
Hermitian solution. If primitivity is deleted the zero vector is
a solution. The source's density retains both corank and
primitivity, so neither example contradicts it.
\end{example}

\subsection{The precise global input}

Let $E/F$ be a CM extension: $F$ is totally real of degree
$f=[F:\mathbb Q]$ and $E$ is a totally imaginary quadratic
extension. Write $\mathcal O_F,\mathcal O_E$ for their rings of
integers, $d_F>0$ for the absolute discriminant of $F$, and
$\mathfrak d_{E/F}$ for the relative discriminant. Assume that
$E/F$ is unramified at every finite prime. For a finite prime $v$
of $F$ put $q_v=\#(\mathcal O_F/v)$ and
$\eta_v=1$ if $v$ splits in $E$, $\eta_v=-1$ if it is inert.
Let $\chi$ be the quadratic Hecke character of $E/F$.
For $j\geq2$ put
\[
                L_j=L_F(j,\chi^j)
                   =\prod_v(1-\eta_v^j q_v^{-j})^{-1}.
\]
Here $\chi^j$ means the trivial character for even $j$ and
$\chi$ for odd $j$. Write $L_1=L_F(1,\chi)>0$ for the value
defined by analytic continuation, or equivalently by the ratio
of the residues of $\zeta_E$ and $\zeta_F$ at $1$.

For a positive definite Hermitian space $V$ over $E$ and a full
$\mathcal O_E$-lattice $L$, its genus is the set of lattices
isometric to $L$ at every finite completion, modulo global
isometry. Its mass is
\[
                \operatorname{Mass}(L)
                   =\sum_{[L']\in\operatorname{gen}(L)}
                                  \frac1{\#\operatorname{Aut}(L')}.
\]
The class set and stabilizers are finite. The class interpretation
and finiteness in the setting used here were proved in
Sections~\ref{al-section} and \ref{acl-section}.
For $n\geq1$, let $M_n$ be the mass of the standard lattice
$\mathcal O_E^n$ with Hermitian form $\overline x^{\,t}y$.

We use the classical Siegel mass formula in its unitary form.
With unimodular local lattices, it states
\begin{equation}\label{md-classical-mass}
 \operatorname{Mass}(L)=
  2d_F^{n^2/2}\,N_{F/\mathbb Q}(\mathfrak d_{E/F})^{n(n+1)/4}
   \left(\prod_{j=1}^n\frac{(j-1)!}{(2\pi)^j}\right)^f
                            \prod_v\beta_{n,v}^{-1}.
\end{equation}
The factor $2$ is the unitary Tamagawa number; the compact
archimedean group has Haar probability measure when interpreting
the class mass. The product has the usual $L(1,\chi)$ convergence
normalization. The formula in this generality is recorded in
Kirschmer, Theorem 4.2.3, from Gan--Yu, \S10.
Terakado--Yu, equation (2.6), gives the same constants.
\cite{kirschmer-herm,terakado-yu-mass}
The latter article imposes an additional dyadic assumption for
its ramified local factors. We do not use those factors: the
unramified local formula needed here has been proved above.
Thus no unramifiedness condition on $F/\mathbb Q$ is being added
to the source hypotheses.

\begin{theorem}\label{md-mass-ratio}
Under the preceding hypotheses, $M_n>0$ and
\[
 M_n=2d_F^{n^2/2}
             \prod_{j=1}^n
                  \left(\frac{(j-1)!}{(2\pi)^j}\right)^f L_j.
\]
In particular, for $n\geq2$,
\begin{equation}\label{md-ratio}
 \frac{M_n}{M_{n-1}}
       =\left((n-1)!\left(\frac{\rho}{2\pi}\right)^n\right)^f
                   d_F^{-1/2}L_n,
             \qquad \rho=d_F^{1/f}.
\end{equation}
Equivalently, with $v=\rho/2$,
\[
 \frac{M_n}{M_{n-1}}
     =\frac{\bigl((n-1)!(v/\pi)^n\bigr)^f}
                  {\sqrt{(2v)^f}}\,L_n.
\]
This is the source mass ratio, including its archimedean
normalization and discriminant power.
\end{theorem}

\begin{proof}
The relative discriminant is the unit ideal.
Substitute \eqref{md-stabilizer-density} in the classical
formula \eqref{md-classical-mass}. Reorganizing the finitely
many rank factors gives the displayed expression for $M_n$.
Equivalently, one first uses the usual $L(1,\chi)$
convergence factor and then the absolutely convergent products
with $j\geq2$; no rearrangement of an unregularized product at
$1$ is needed. Positivity also follows directly because the
genus contains the standard lattice and all summands are positive.

On division by the expression in rank $n-1$, the factor $2$
and all factors with $j<n$ cancel. The remaining discriminant
power is $d_F^{(2n-1)/2}$. Since $\rho^f=d_F$, this is
$\rho^{nf}d_F^{-1/2}$, which proves \eqref{md-ratio}.
Replacing $\rho$ by $2v$ gives its last form.
\end{proof}

\subsection{The source field and Euler-product conventions}

The source initially specifies a totally imaginary number field $E$
with an involution whose image under every complex embedding is complex
conjugation. Its fixed field is indeed a totally real field $F$ with
$[E:F]=2$. The involution cannot be trivial, since that would make
every embedding real. Choose $\alpha$ not fixed by it. For $x\in E$,
the element
\[
 b=\frac{x-\overline x}{\alpha-\overline\alpha}
\]
is fixed, and $x-b\alpha$ is fixed as well. Thus $E=F+F\alpha$,
and the sum has dimension two. Every embedding of $F$ extends to
an embedding of $E$; its values on $F$ are real by the specified
conjugation identity. This proves the asserted CM interpretation.

\begin{proposition}\label{md-source-conventions}
Let complex conjugation on $E$ be denoted by a bar.
The following condition is equivalent to finite unramifiedness
of $E/F$:
for every maximal ideal $P\subset\mathcal O_E$ there exists
$\theta\in\mathcal O_E$ with $\theta-\overline\theta\notin P$.
Under this condition,
\[
 [E:\mathbb Q]=2f,\qquad |d_E|=d_F^2,
 \qquad |d_E|^{1/(2f)}=\rho.
\]
The product indexed by primes $P$ of $E$ with local factor
\[
 \begin{cases}
 (1-(-1)^n(NP)^{-n/2})^{-1},&P=\overline P,\\
 (1-(NP)^{-n})^{-1/2},&P\ne\overline P
 \end{cases}
\]
equals $L_n$ for $n\geq2$, where the positive square root is used.
\end{proposition}

\begin{proof}
At a ramified prime of a quadratic extension, conjugation acts
trivially on the residue field, so every integral
$\theta-\overline\theta$ belongs to $P$.
At an inert unramified prime, conjugation is the nontrivial
automorphism of its quadratic residue-field extension. Choose
a residue outside the base field and lift it to $\mathcal O_E$.
At a split prime, $P$ and $\overline P$ are distinct, and the
Chinese remainder theorem gives $\theta\equiv0\pmod P$ and
$\theta\equiv1\pmod{\overline P}$. In particular
$\overline\theta\equiv1\pmod P$. This proves both implications.

The discriminant formula in a tower gives
$|d_E|=d_F^2N(\mathfrak d_{E/F})=d_F^2$, and the degree
formula gives $[E:\mathbb Q]=2f$.
At an inert prime $v$, there is one $P$ and $NP=q_v^2$.
Its displayed factor is $(1-(-1)^nq_v^{-n})^{-1}$.
At a split prime there are two primes of norm $q_v$, and
their positive square-root factors multiply to
$(1-q_v^{-n})^{-1}$. These are exactly the $v$-factors of $L_n$.
Absolute convergence for $n\geq2$ justifies grouping them.
\end{proof}

\begin{corollary}\label{md-euler-bound}
For every integer $n\geq2$,
\[
                         |\log L_n|\leq f\log\zeta(n).
\]
Thus the recurrence and Euler-factor bound assumed in
Section~\ref{as-section} follow from the stated classical
mass formula with the actual source hypotheses.
\end{corollary}

\begin{proof}
For $0<x<1$ and $\epsilon\in\{1,-1\}$,
$|\log(1-\epsilon x)|\leq-\log(1-x)$; this follows from
$1+x\leq(1-x)^{-1}$ and monotonicity of the logarithm.
Apply this to every local factor to obtain
$|\log L_n|\leq\log\zeta_F(n)$.
Above a rational prime $p$, there are at most $f$ prime ideals
of $F$, each of norm at least $p$. Therefore
\[
 \prod_{v\mid p}(1-q_v^{-n})^{-1}
                             \leq(1-p^{-n})^{-f}.
\]
Multiplying the absolutely convergent products proves
$\zeta_F(n)\leq\zeta(n)^f$ and the assertion.
\end{proof}

The finite-field formula proves the source primitive-frame density
including all Gram coranks. The mass ratio removes the mass input
from the archimedean size calculation, subject only to the classical
global mass theorem explicitly identified above. The source-space
dimension and formal-degree inputs of that calculation are distinct
representation-theoretic assertions and are not inferred from the
mass ratio.


% Exact column dimensions and finite-group fixed spaces.
% Labels in this fragment use the prefix cd-.
